\documentclass[11pt,letterpaper]{article}

\usepackage[spanish,es-noshorthands]{babel}   % '<' y '>' se usan en lasProducciones
\usepackage{fontspec}
\setmainfont[Path=fonts/,UprightFont=NotoSans-Regular.ttf,BoldFont=NotoSans-Bold.ttf,ItalicFont=NotoSans-Regular.ttf,ItalicFeatures={FakeSlant=0.2}]{NotoSans-Regular.ttf}
\usepackage{microtype}
\usepackage{amsmath}
\usepackage[letterpaper,margin=2.15cm,headheight=14pt]{geometry}
\usepackage{graphicx}
\usepackage{caption}
\usepackage{booktabs}
\usepackage{tabularx}
\usepackage{longtable}
\usepackage{pdflscape}
\usepackage{array}
\usepackage[table]{xcolor}
\usepackage{enumitem}
\usepackage{fancyhdr}
\usepackage{hyperref}
\usepackage{titlesec}
\usepackage{needspace}
\usepackage{etoolbox}

\definecolor{azul}{HTML}{2B5F9E}
\definecolor{azulclaro}{HTML}{DCEAF7}
\definecolor{grisclaro}{HTML}{F3F5F6}
\definecolor{verde}{HTML}{1E6B3A}
\definecolor{rojo}{HTML}{9B2226}

\hypersetup{
  pdftitle={Informe de Laboratorio 03: Diseno y Validacion de Gramaticas},
  pdfauthor={Juan Muñoz Veloso; Fernando Valdés Gómez},
  pdfsubject={Teoría de la Computación INFO1148},
  colorlinks=true,
  linkcolor=azul,
  urlcolor=azul
}

\setlength{\emergencystretch}{3em}
\setlength{\parindent}{0pt}
\setlength{\parskip}{0.4em}
\setlength{\tabcolsep}{7pt}
\renewcommand{\arraystretch}{1.18}
\setlist[itemize]{leftmargin=1.35em,itemsep=0.15em,topsep=0.2em}
\newcolumntype{Y}{>{\raggedright\arraybackslash}X}
\newcolumntype{L}[1]{>{\raggedright\arraybackslash}p{#1}}
\titleformat{\section}{\Large\bfseries\color{azul}}{}{0pt}{}
\titleformat{\subsection}{\large\bfseries\color{azul}}{}{0pt}{}

\pagestyle{fancy}
\fancyhf{}
\fancyfoot[L]{\small INFO1148 · Laboratorio 3}
\fancyfoot[R]{\thepage}
\renewcommand{\headrulewidth}{0pt}
\renewcommand{\footrulewidth}{0pt}

\captionsetup{
  font=normalsize,
  labelfont=bf,
  labelsep=period,
  justification=raggedright,
  singlelinecheck=false,
  hypcap=false
}
\renewcommand{\figurename}{Figura}

\newcommand{\apartado}[1]{\par\Needspace{4\baselineskip}\medskip\noindent\textbf{#1}\par\nobreak\smallskip}

% Notacion de la gramatica: los no terminales van entre angle brackets.
\newcommand{\nt}[1]{\texttt{\textless #1\textgreater}}
\newcommand{\tm}[1]{\texttt{#1}}
% Regla de produccion dentro de una tabla.
\newcommand{\regla}[2]{\nt{#1} $\rightarrow$ \texttt{#2}}

% Cada figura se mantiene unida a su titulo, leyenda y descripcion.
\newcommand{\figuraaccesible}[5]{%
  \par\medskip
  \noindent\begin{minipage}{\textwidth}
    {\large\bfseries\color{azul}#2\par}
    \begin{center}
      \includegraphics[width=0.90\textwidth,height=0.40\textheight,keepaspectratio]{#1}
      \captionof{figure}{#3}\label{#5}
    \end{center}
    #4\par
  \end{minipage}\par\medskip
}

% El arbol de la extension se muestra en una pagina apaisada.
\newcommand{\figuraamplia}[5]{%
  \clearpage
  \begin{landscape}
  {\large\bfseries\color{azul}#2\par}
  \begin{center}
    \includegraphics[width=0.98\linewidth,height=0.76\textheight,keepaspectratio]{#1}
    \captionof{figure}{#3}\label{#5}
  \end{center}
  #4\par
  \end{landscape}
}

\pretocmd{\subsection}{\Needspace{9\baselineskip}}{}{}

\begin{document}
\pagenumbering{arabic}

\begin{center}
  {\huge\bfseries INFORME DE LABORATORIO \#3\par}
  \vspace{0.2cm}
  {\Large\bfseries Diseño y Validación de Gramáticas\par}
  \vspace{0.55cm}
\end{center}

\begin{tabularx}{\textwidth}{|L{0.38\textwidth}|Y|}
\hline
\rowcolor{azulclaro}\textbf{Antecedente} & \textbf{Información}\\\hline
Asignatura & Teoría de la Computación INFO1148\\\hline
Carrera & Ingeniería Civil en Informática\\\hline
Docente & Marcos Lévano\\\hline
Integrantes & Juan Muñoz Veloso y Fernando Valdés Gómez\\\hline
Fecha & 20 de octubre de 2026\\\hline
\end{tabularx}

\section*{1. Introducción}
En este laboratorio se diseña, analiza y valida una Gramática Libre de Contexto (GLC) capaz de
representar la estructura sintáctica del lenguaje de programación simplificado del enunciado. Se
identifican los terminales y los no terminales del lenguaje, se define formalmente la gramática
$G=(V,\Sigma,P,S)$, se construyen sus producciones, se realizan derivaciones por la izquierda, se
construye el árbol sintáctico de las cadenas solicitadas y se valida el comportamiento de la
gramática frente a ocho casos, cuatro correctos y cuatro incorrectos. Como extensión creativa se
incorpora la construcción de repetición (\texttt{mientras ... fin\_mientras}) al lenguaje, lo que
obliga a ampliar tanto el conjunto de terminales como el de no terminales. Toda la validación se
apoya en un analizador descendente recursivo con retroceso entregado junto al informe. Además, las
tres derivaciones principales se reprodujeron en JFLAP 7.1: el archivo
\texttt{gramatica/\allowbreak GLC\_Lab03.cfg} usa una codificación isomorfa de un carácter por
símbolo, explicada en el anexo D, y las capturas de los árboles se incluyen como evidencia. La
lectura del archivo se comprobó con la propia biblioteca de JFLAP 7.1; el detalle está en el
anexo D.

\clearpage
\section*{2. Desarrollo}

\subsection*{2.1 Pregunta 1 -- Identificación formal de la gramática}

\apartado{a) Terminales y su clasificación}

Los terminales son los símbolos que no admiten ninguna derivación: son las fichas concretas del
lenguaje. Se clasifican según su función en el lenguaje.

\begin{longtable}{|L{3.3cm}|L{5.1cm}|L{\dimexpr\textwidth-3.3cm-5.1cm-6\tabcolsep-4\arrayrulewidth\relax}|}
\hline
\rowcolor{azulclaro}\textbf{Categoría} & \textbf{Terminales} & \textbf{Observación}\\\hline
\endfirsthead
\hline
\rowcolor{azulclaro}\textbf{Categoría} & \textbf{Terminales} & \textbf{Observación}\\\hline
\endhead

Palabras reservadas & \texttt{Begin}, \texttt{End}, \texttt{Programa}, \texttt{salida} &
No admiten otra función: el analizador léxico los reconoce como palabras completas.\\\hline
Operadores aritméticos & \texttt{+}, \texttt{-}, \texttt{*}, \texttt{/}, \texttt{\%} &
Binarios. Cada expresión admite una sola operación, sin reglas de precedencia.\\\hline
Operador de asignación & \texttt{=} &
Único en el lenguaje; separa el identificador del valor.\\\hline
Símbolos especiales & \texttt{\#} &
Termina asignaciones y salidas. El lazo se cierra con \texttt{fin\_mientras}.\\\hline
Identificadores & Clase léxica \texttt{IDENT} &
Forma: \texttt{[A-Za-z]\allowbreak[A-Za-z0-9\_]*}. El guion bajo pertenece a la clase, no es un terminal aparte.\\\hline
Números & Clase léxica \texttt{NUM} &
Forma: \texttt{[0-9]+(\allowbreak\textbackslash.[0-9]+)?}. Cubre \texttt{20}, \texttt{0} y \texttt{1.345}.\\\hline
\end{longtable}

La extensión de la sección 3 agrega dos palabras reservadas (\texttt{mientras},
\texttt{fin\_mientras}) y cuatro operadores relacionales (\texttt{<}, \texttt{>}, \texttt{\textless=},
\texttt{\textgreater=}), que se describen en detalle más adelante.

\apartado{b) No terminales y función de cada uno}

{\raggedright
Los no terminales identificados son \nt{programa}, \nt{lista\_instrucciones},
\nt{instruccion}, \nt{asignacion}, \nt{asignacion\_simple}, \nt{asignacion\_expresion},
\nt{valor}, \nt{expresion}, \nt{operando}, \nt{operador} y \nt{salida\_instr}.
La extensión agrega \nt{mientras\_instr}, \nt{condicion} y \nt{operador\_rel}.
La función de cada uno se explica en el apartado 2.2 b).\par}

\apartado{c) Símbolo inicial}

El símbolo inicial es \nt{programa}. Toda cadena del lenguaje completo se obtiene derivando a
partir de ese no terminal, que es el único que aparece en el encabezado
\texttt{Begin Programa <id> ... End}. Las cadenas que son sólo un fragmento del programa (una
instrucción suelta) se derivan a partir de \nt{instruccion} o de \nt{lista\_instrucciones}, que son
subgramáticas de la principal.

\apartado{d) Definición formal $G=(V,\Sigma,P,S)$}

Una GLC es un cuádruplo $G=(V,\Sigma,P,S)$: $V$ contiene los no terminales, $\Sigma$ los terminales,
$P$ las producciones y $S$ el símbolo inicial. La definición siguiente incluye la extensión.
La gramática base tiene 11 no terminales, 13 terminales y 20 producciones. La versión extendida
agrega tres no terminales, seis terminales y siete producciones, para un total de 14, 19 y 27.

\begin{longtable}{|L{2.7cm}|L{\dimexpr\textwidth-2.7cm-4\tabcolsep-3\arrayrulewidth\relax}|}
\hline
\rowcolor{azulclaro}\textbf{Componente} & \textbf{Definición}\\\hline
\endfirsthead
\hline
\rowcolor{azulclaro}\textbf{Componente} & \textbf{Definición}\\\hline
\endhead

{\small $V$} & $\{\,$programa, lista\_instrucciones, instruccion, asignacion, asignacion\_simple,
asignacion\_expresion, valor, expresion, operando, operador, salida\_instr, mientras\_instr,
condicion, operador\_rel$\}$, con $|\,V\,|=14$.\\\hline
{\small $\Sigma$} & $\{\,$Begin, End, Programa, salida, mientras, fin\_mientras, \texttt{+}, \texttt{-}, \allowbreak
\texttt{*}, \texttt{/}, \texttt{\%}, \texttt{=}, \texttt{\#}, \texttt{<}, \texttt{>}, \allowbreak
\texttt{\textless=}, \texttt{\textgreater=}, IDENT, NUM$\}$, con $|\,\Sigma\,|=19$.\\\hline
$P$ & Conjunto de 27 reglas de producción, numeradas de P01 a P27 (sección 2.2).\\\hline
$S$ & \nt{programa}.\\\hline
\end{longtable}

Los conjuntos $V$ y $\Sigma$ son disjuntos y todas las reglas tienen la forma $X\rightarrow\alpha$ con
$X\in V$, que es la condición de gramática libre de contexto.

\apartado{e) Por qué no enumerar identificadores y números}

Los identificadores y los números admiten infinitos valores. Enumerarlos exigiría infinitas
producciones, mientras que una GLC debe tener un conjunto finito de reglas. Se reconocen por su
forma durante el análisis léxico; la gramática sintáctica recibe las clases \texttt{IDENT} y
\texttt{NUM}. También sería válido describirlos con reglas recursivas a nivel de caracteres,
pero aquí se separan ambos niveles para mantener legibles las producciones.

La solución adoptada es la habitual en este tipo de lenguajes: los identificadores y los números se
tratan como dos clases léxicas, \texttt{IDENT} y \texttt{NUM}, descritas por las expresiones
regulares \texttt{[A-Za-z]\allowbreak[A-Za-z0-9\_]*} y \texttt{[0-9]+(\allowbreak\textbackslash.[0-9]+)?}. La gramática
trabaja con esas clases como si fueran fichas atómicas. En el anexo A se muestra cómo podría
describirse la misma gramática en su versión carácter a carácter, con reglas del tipo
\nt{identificador} $\rightarrow$ \nt{letra} \nt{resto\_id}, lo que confirma que la decisión es una
simplificación legítima y no una omisión.

\subsection*{2.2 Pregunta 2 -- Construcción de la Gramática Libre de Contexto}

\apartado{a) Producciones}

La tabla siguiente contiene el conjunto completo de reglas $P$, con la justificación de cada una.
En el lado derecho se omiten los ángulos para abreviar: los nombres definidos en $V$ son
no terminales; los restantes símbolos pertenecen a $\Sigma$.
Las reglas P01 a P21, excepto P06, describen el lenguaje del enunciado; P06 y P22 a P27 corresponden a
la extensión creativa de la sección 3 y se analizan allí.

\begin{longtable}{|L{1.2cm}|L{\dimexpr\textwidth-1.2cm-5.0cm-6\tabcolsep-4\arrayrulewidth\relax}|L{5.0cm}|}
\hline
\rowcolor{azulclaro}\textbf{Regla} & \textbf{Producción} & \textbf{Justificación}\\\hline
\endfirsthead
\hline
\rowcolor{azulclaro}\textbf{Regla} & \textbf{Producción} & \textbf{Justificación}\\\hline
\endhead

P01 & \regla{programa}{Begin Programa IDENT lista\_instrucciones End} &
Encabezado del programa: la palabra \texttt{Programa} precede al nombre del programa y el cuerpo
va entre \texttt{Begin} y \texttt{End}.\\\hline
P02 & \regla{lista\_instrucciones}{instruccion} &
Caso base: la lista puede tener una sola instrucción.\\\hline
P03 & \regla{lista\_instrucciones}{instruccion lista\_instrucciones} &
Caso recursivo: permite tantas instrucciones como se necesiten.\\\hline
P04 & \regla{instruccion}{asignacion} &
Una instrucción puede ser una asignación.\\\hline
P05 & \regla{instruccion}{salida\_instr} &
Una instrucción puede ser una salida por pantalla.\\\hline
P06 & \regla{instruccion}{mientras\_instr} &
Extensión creativa: lazo repetitivo.\\\hline
P07 & \regla{asignacion}{asignacion\_simple} &
Distingue la asignación de un valor simple.\\\hline
P08 & \regla{asignacion}{asignacion\_expresion} &
Distingue la asignación de un resultado aritmético.\\\hline
P09 & \regla{asignacion\_simple}{IDENT = valor \#} &
Asignación simple: identificador, igual, valor y terminador.\\\hline
P10 & \regla{asignacion\_expresion}{IDENT = expresion \#} &
Asignación con operación aritmética.\\\hline
P11 & \regla{valor}{NUM} &
Un valor puede ser un número literal.\\\hline
P12 & \regla{valor}{IDENT} &
Un valor puede ser el contenido de otra variable (\texttt{temporal = x \#}).\\\hline
P13 & \regla{expresion}{operando operador operando} &
Expresión aritmética binaria, la forma que usa el programa de referencia.\\\hline
P14 & \regla{operando}{IDENT} &
Un operando puede ser una variable.\\\hline
P15 & \regla{operando}{NUM} &
Un operando puede ser una constante numérica.\\\hline
P16 & \regla{operador}{+} & Operador aritmético de suma.\\\hline
P17 & \regla{operador}{-} & Operador aritmético de resta.\\\hline
P18 & \regla{operador}{*} & Operador aritmético de multiplicación.\\\hline
P19 & \regla{operador}{/} & Operador aritmético de división.\\\hline
P20 & \regla{operador}{\%} & Operador aritmético de módulo.\\\hline
P21 & \regla{salida\_instr}{salida IDENT \#} &
Instrucción de salida: palabra reservada, identificador y terminador.\\\hline
P22 & \regla{mientras\_instr}{mientras condicion lista\_instrucciones fin\_mientras} &
Extensión creativa: cabecera, condición, cuerpo y cierre del lazo.\\\hline
P23 & \regla{condicion}{operando operador\_rel operando} &
La condición compara dos operandos.\\\hline
P24 & \regla{operador\_rel}{<} & Operador relacional menor que.\\\hline
P25 & \regla{operador\_rel}{>} & Operador relacional mayor que.\\\hline
P26 & \regla{operador\_rel}{\textless=} & Operador relacional menor o igual que.\\\hline
P27 & \regla{operador\_rel}{\textgreater=} & Operador relacional mayor o igual que.\\\hline
\end{longtable}

\apartado{b) No terminales y función de cada uno}

Los no terminales son los símbolos que se pueden derivar: representan las estructuras del lenguaje.
La tabla siguiente justifica el papel de cada uno dentro de la gramática.

\begin{longtable}{|L{4.8cm}|L{\dimexpr\textwidth-4.8cm-4\tabcolsep-3\arrayrulewidth\relax}|}
\hline
\rowcolor{azulclaro}\textbf{No terminal} & \textbf{Función}\\\hline
\endfirsthead
\hline
\rowcolor{azulclaro}\textbf{No terminal} & \textbf{Función}\\\hline
\endhead

\nt{programa} & Encabezado y cuerpo completo del programa: \texttt{Begin Programa <id> ... End}.\\\hline
\nt{lista\_instrucciones} & Secuencia de una o más instrucciones del cuerpo del programa.\\\hline
\nt{instruccion} & Una instrucción del lenguaje: asignación, salida o lazo.\\\hline
\nt{asignacion} & Encabezamiento común de los dos tipos de asignación.\\\hline
\nt{asignacion\_simple} & Asignación de un valor simple: \texttt{<id> = <valor> \#}.\\\hline
\nt{asignacion\_expresion} & Asignación del resultado de una operación: \texttt{<id> = <expresion> \#}.\\\hline
\nt{valor} & Valor simple asignable: un número o un identificador.\\\hline
\nt{expresion} & Expresión aritmética binaria: \texttt{<operando> <operador> <operando>}.\\\hline
\nt{operando} & Cada uno de los dos operandos de una expresión.\\\hline
\nt{operador} & Operador aritmético binario: \texttt{+ - * / \%}.\\\hline
\nt{salida\_instr} & Instrucción de salida por pantalla: \texttt{salida <id> \#}.\\\hline
\nt{mientras\_instr} & Extensión creativa: lazo \texttt{mientras <condicion> ... fin\_mientras}.\\\hline
\nt{condicion} & Extensión creativa: comparación entre dos operandos.\\\hline
\nt{operador\_rel} & Extensión creativa: operador relacional \texttt{< > \textless= \textgreater=}.\\\hline
\end{longtable}

Dos nombres llevan un sufijo (\nt{salida\_instr}, \nt{mientras\_instr}) porque la terminal
\texttt{salida} y la terminal \texttt{mientras} existen en $\Sigma$. La condición formal
$V\cap\Sigma=\emptyset$ exige que un no terminal y un terminal no se llamen igual, de modo que se
distinguen los símbolos \nt{salida\_instr} y \nt{mientras\_instr} de las palabras reservadas
correspondientes.

\apartado{c) Diferenciación entre asignación simple, asignación con expresión y salida}

El enunciado pide diferenciar explícitamente las tres construcciones. La diferencia no es de
nombre, sino de posición dentro de la regla y de la forma del lado derecho:

\begin{itemize}
  \item \textbf{Asignación simple (P09):} \nt{asignacion\_simple} $\rightarrow$ \texttt{IDENT = valor \#}.
  El valor es un único operando simple, un número o un identificador. Ejemplos:
  \texttt{x = 20 \#}, \texttt{temporal = x \#}.
  \item \textbf{Asignación con expresión aritmética (P10):}
  \nt{asignacion\_expresion} $\rightarrow$ \texttt{IDENT = expresion \#}, donde
  \nt{expresion} $\rightarrow$ \nt{operando} \nt{operador} \nt{operando} exige obligatoriamente un
  operador binario. Ejemplos: \texttt{z = x + y \#}, \texttt{k = h * y \#}.
  \item \textbf{Instrucción de salida (P21):} \nt{salida\_instr} $\rightarrow$ \texttt{salida IDENT \#}.
  No hay signo igual ni expresión: el argumento es directamente un identificador. Ejemplo:
  \texttt{salida z \#}.
\end{itemize}

La asignación simple y la asignación con expresión se distinguen porque la segunda exige un
operador binario. En los casos ensayados, cada cadena aceptada tiene un único árbol sintáctico.
Este conteo comprueba los ejemplos, pero por sí solo no demuestra la ausencia de ambigüedad
para todas las cadenas del lenguaje.

\apartado{d) Verificación con el programa de referencia}

El programa de referencia del enunciado se reconoce completo: el analizador lo acepta como
derivado de \nt{programa}. La derivación por la izquierda tiene 108 pasos; las formas sentenciales
comienzan así:

\begin{quote}\small\ttfamily
\nt{programa}\\
$\Rightarrow$ Begin Programa operativo\_01\_02 \nt{lista\_instrucciones} End\\
$\Rightarrow$ Begin Programa operativo\_01\_02 \nt{instruccion} \nt{lista\_instrucciones} End\\
$\Rightarrow$ Begin Programa operativo\_01\_02 \nt{asignacion} \nt{lista\_instrucciones} End\\
$\Rightarrow$ Begin Programa operativo\_01\_02 \nt{asignacion\_simple} \nt{lista\_instrucciones} End\\
$\Rightarrow$ Begin Programa operativo\_01\_02 x = \nt{valor} \# \nt{lista\_instrucciones} End\\
$\Rightarrow$ Begin Programa operativo\_01\_02 x = 1.345 \# \nt{lista\_instrucciones} End\\
$\Rightarrow$ Begin Programa operativo\_01\_02 x = 1.345 \# \nt{instruccion} \nt{lista\_instrucciones} End\\
$\Rightarrow$ $\ldots$
\end{quote}

La derivación completa está en el archivo indicado en el anexo B y en
el archivo \texttt{evidencias/\allowbreak evidencias\_lab03.txt}. El programa contiene 91 fichas
(\texttt{token}), contando las palabras reservadas, los identificadores, los números, los
operadores y los terminadores.

\subsection*{2.3 Pregunta 3 -- Derivaciones y árbol sintáctico}

\apartado{a) Derivación por la izquierda de \texttt{z = x + y \#}}

La cadena pedida es una instrucción, por lo que la derivación se inicia en \nt{instruccion}, el no
terminal que la genera dentro de \nt{lista\_instrucciones}. Cada paso reemplaza el no terminal más
a la izquierda de la forma sentencial, que es lo que caracteriza a una derivación por la
izquierda. Se usa notación con lexemas: al aplicar una regla con \texttt{IDENT} o \texttt{NUM},
se muestra directamente el texto del token reconocido. Así, \texttt{z}, \texttt{x} e \texttt{y}
son ocurrencias de \texttt{IDENT}; no son nuevos terminales ni requieren reglas adicionales.
Los árboles siguen la misma convención en sus hojas.

\Needspace{14\baselineskip}
\begin{quote}\small\ttfamily
\nt{instruccion}\\
$\Rightarrow$ \nt{asignacion} \quad [P04]\\
$\Rightarrow$ \nt{asignacion\_expresion} \quad [P08]\\
$\Rightarrow$ z = \nt{expresion} \# \quad [P10]\\
$\Rightarrow$ z = \nt{operando} \nt{operador} \nt{operando} \# \quad [P13]\\
$\Rightarrow$ z = x \nt{operador} \nt{operando} \# \quad [P14]\\
$\Rightarrow$ z = x + \nt{operando} \# \quad [P16]\\
$\Rightarrow$ z = x + y \# \quad [P14]
\end{quote}

\Needspace{8\baselineskip}
\Needspace{12\baselineskip}
\apartado{b) Derivación por la izquierda de \texttt{salida z \#}}

\begin{quote}\small\ttfamily
\nt{instruccion}\\
$\Rightarrow$ \nt{salida\_instr} \quad [P05]\\
$\Rightarrow$ salida z \# \quad [P21]
\end{quote}

La derivación es más corta porque la instrucción de salida no tiene partes que expandir: los tres
símbolos de la derecha de P21 son ya terminales. El árbol adicional de esta cadena se
incluye en el anexo E.

\figuraaccesible{figuras/arbol_z_xy.pdf}{c) Árbol sintáctico de \texttt{z = x + y \#}}{Árbol de
derivación de la cadena \texttt{z = x + y \#} obtenido con la herramienta de validación entregada
junto al informe.}{La raíz es \nt{instruccion}; los nodos internos son los no terminales que aún
debían expandirse y las hojas son los terminales leídos, de izquierda a derecha: \texttt{z},
\texttt{=}, \texttt{x}, \texttt{+}, \texttt{y} y \texttt{\#}.}{fig:arbol-zxy}

\apartado{d) Relación entre la raíz, los no terminales, los terminales y las hojas}

La raíz \nt{instruccion} identifica la estructura que se deriva. Los nodos internos son
no terminales y sus hijos representan el lado derecho de la producción aplicada. Las hojas
son terminales; para \texttt{IDENT} y \texttt{NUM} se muestra el lexema del token.
Leer las hojas de izquierda a derecha reconstruye \texttt{z = x + y \#}.
El árbol conserva la jerarquía de la derivación, mientras que los pasos por la izquierda
indican en qué orden se expandieron los no terminales.

\Needspace{14\baselineskip}
\apartado{e) Validación de las tres cadenas adicionales}

Las tres cadenas de la Pregunta 3.14 se generan con la gramática propuesta:

\Needspace{8\baselineskip}
\begin{longtable}{|L{4.0cm}|L{2.6cm}|L{\dimexpr\textwidth-4.0cm-2.6cm-6\tabcolsep-4\arrayrulewidth\relax}|}
\hline
\rowcolor{azulclaro}\textbf{Cadena} & \textbf{¿Generable?} & \textbf{Justificación}\\\hline
\endfirsthead
\hline
\rowcolor{azulclaro}\textbf{Cadena} & \textbf{¿Generable?} & \textbf{Justificación}\\\hline
\endhead

\texttt{x = 20 \#} & Sí & Asignación simple (P09): \nt{valor} $\rightarrow$ \texttt{NUM}, cuyo lexema es \texttt{20}.\\\hline
\texttt{z = x * y \#} & Sí & Asignación con expresión (P10): \texttt{*} se deriva con P18 y ambos operandos son identificadores.\\\hline
\texttt{salida temporal \#} & Sí & Instrucción de salida (P21) con el identificador \texttt{temporal}.\\\hline
\end{longtable}

\subsection*{2.4 Pregunta 4 -- Validación y detección de errores}

\apartado{a) Método de validación}

La validación se realizó con un analizador descendente recursivo con retroceso completo, entregado
en el archivo \texttt{herramientas/\allowbreak validar\_gramatica.py}. El analizador tokeniza el código fuente
(ignorando los comentarios \texttt{/* ... */}), aplica exactamente las 27 reglas de producción
entregadas y devuelve, cuando la cadena es válida, el árbol sintáctico correspondiente. El motor de
análisis de JFLAP 7.1 trabaja carácter a carácter: las letras mayúsculas son variables y cualquier
otro carácter es terminal. Para que el archivo \texttt{gramatica/\allowbreak GLC\_Lab03.cfg}
represente la misma gramática, cada símbolo se codifica con un solo carácter (anexo D). El
archivo se cargó con la propia biblioteca de JFLAP 7.1 y reporta 14 variables, 19 terminales y 27
producciones, con \texttt{S} como símbolo inicial (anexo D). Las capturas del anexo D muestran
las derivaciones de las tres cadenas principales en JFLAP.

\apartado{b) Resultado del análisis de los ocho casos}

\begin{longtable}{|L{1.3cm}|L{4.3cm}|L{2.1cm}|L{\dimexpr\textwidth-1.3cm-4.3cm-2.1cm-8\tabcolsep-5\arrayrulewidth\relax}|}
\hline
\rowcolor{azulclaro}\textbf{Caso} & \textbf{Cadena} & \textbf{Resultado} & \textbf{Análisis}\\\hline
\endfirsthead
\hline
\rowcolor{azulclaro}\textbf{Caso} & \textbf{Cadena} & \textbf{Resultado} & \textbf{Análisis}\\\hline
\endhead

1 & \texttt{x = 10 \#} & \textcolor{verde}{Aceptada} &
Asignación simple: P09 con \nt{valor} $\rightarrow$ \texttt{NUM}.\\\hline
2 & \texttt{z=x+y\#} & \textcolor{verde}{Aceptada} &
Asignación con expresión: P10, P13 y P16.\\\hline
3 & \texttt{salida z \#} & \textcolor{verde}{Aceptada} &
Instrucción de salida: P05 y P21.\\\hline
4 & \texttt{x 10 \#} & \textcolor{rojo}{Rechazada} &
Falta el operador de asignación \texttt{=}: la regla P09 exige la secuencia
\texttt{IDENT = valor \#}.\\\hline
5 & \texttt{z=+xy\#} & \textcolor{rojo}{Rechazada} &
El signo \texttt{+} queda donde la gramática espera un operando: P13 exige
\nt{operando} \nt{operador} \nt{operando} y el lenguaje no admite signos unarios.\\\hline
6 & \texttt{salida \#} & \textcolor{rojo}{Rechazada} &
Falta el argumento de la instrucción: P21 exige \texttt{salida IDENT \#}.\\\hline
7 & \texttt{x = 20} & \textcolor{rojo}{Rechazada} &
Falta el terminador de instrucción \texttt{\#}.\\\hline
8 & \texttt{Begin Programa prueba x = 5 \# salida x \# End} & \textcolor{verde}{Aceptada} &
Programa completo: P01 con dos instrucciones en el cuerpo.\\\hline
\end{longtable}

\apartado{c) Alcance de los rechazos}

El caso 5 tiene un error sintáctico, no léxico: \texttt{xy} sí es un identificador válido,
pero \texttt{+} no puede aparecer antes del primer operando. En los casos 4, 6 y 7 faltan,
respectivamente, \texttt{=}, el identificador de salida y \texttt{\#}. El análisis comprueba
la estructura; no comprueba declaración de variables, tipos ni divisiones por cero.

\section*{3. Extensión de la gramática}

De las cinco opciones propuestas se eligió la \textbf{opción B, repetición}, porque introduce una
construcción que anida otras instrucciones y obliga a ampliar la gramática en los dos sentidos:
nuevos terminales y nuevos no terminales.

\apartado{a) Sintaxis de la construcción}

\begin{quote}
\texttt{mientras} \textit{comparación} \textit{instrucciones} \texttt{fin\_mientras}
\end{quote}

La comparación sigue el patrón \nt{operando} \nt{operador\_rel} \nt{operando}, con los operadores
relacionales \texttt{<}, \texttt{>}, \texttt{\textless=} y \texttt{\textgreater=}. El cuerpo del
lazo es una lista de instrucciones con la misma sintaxis del programa, por lo que puede contener
asignaciones, salidas e incluso otros lazos.

\apartado{b) Terminales y no terminales nuevos}

Se agregan a $\Sigma$ los terminales \texttt{mientras} y \texttt{fin\_mientras} (palabras
reservadas) y \texttt{<}, \texttt{>}, \texttt{\textless=}, \texttt{\textgreater=} (operadores
relacionales). Se agregan a $V$ los no terminales \nt{mientras\_instr}, \nt{condicion} y
\nt{operador\_rel}. Las reglas nuevas son:

\Needspace{10\baselineskip}
\begin{longtable}{|L{1.1cm}|L{\dimexpr\textwidth-1.1cm-4\tabcolsep-3\arrayrulewidth\relax}|}
\hline
\rowcolor{azulclaro}\textbf{Regla} & \textbf{Producción}\\\hline
\endfirsthead
\hline
\rowcolor{azulclaro}\textbf{Regla} & \textbf{Producción}\\\hline
\endhead

P06 & \regla{instruccion}{mientras\_instr}\\\hline
P22 & \regla{mientras\_instr}{mientras condicion lista\_instrucciones fin\_mientras}\\\hline
P23 & \regla{condicion}{operando operador\_rel operando}\\\hline
P24 & \regla{operador\_rel}{<}\\\hline
P25 & \regla{operador\_rel}{>}\\\hline
P26 & \regla{operador\_rel}{\textless=}\\\hline
P27 & \regla{operador\_rel}{\textgreater=}\\\hline
\end{longtable}

Con P06 el lazo pasa a ser una instrucción más y, por tanto, puede aparecer en el cuerpo del
programa y también dentro del cuerpo de otro lazo. La condición reutiliza \nt{operando}, lo que
permite comparar variables y constantes numéricas. No admite expresiones aritméticas como operandos de la comparación.

\Needspace{12\baselineskip}
\apartado{c) Cadena válida y cadena inválida}

\begin{longtable}{|L{2.3cm}|L{\dimexpr\textwidth-2.3cm-4\tabcolsep-3\arrayrulewidth\relax}|}
\hline
\rowcolor{azulclaro}\textbf{Caso} & \textbf{Cadena}\\\hline
\endfirsthead
\hline
\rowcolor{azulclaro}\textbf{Caso} & \textbf{Cadena}\\\hline
\endhead

\textcolor{verde}{Válida} &
\texttt{Begin Programa lazo mientras x < 10 x = x + 1 \# fin\_mientras salida x \# End}\\\hline
\textcolor{rojo}{Inválida} &
\texttt{Begin Programa lazo mientras x < 10 x = x + 1 \# salida x \# End}\\\hline
\end{longtable}

La cadena inválida omite la palabra reservada \texttt{fin\_mientras}: el lazo queda abierto y el
analizador rechaza la cadena, porque P22 exige que al cuerpo del lazo le siga \texttt{fin\_mientras}.
Una segunda forma de error sería un operador relacional ausente, como en
\texttt{mientras x 10 ... fin\_mientras}.

\Needspace{30\baselineskip}
\apartado{d) Derivación de ejemplo}

La derivación por la izquierda de la instrucción de lazo se obtiene del mismo modo que las
anteriores:

Se usan estas abreviaturas tipográficas de los mismos no terminales:

$L$: \nt{lista\_instrucciones}; $I$: \nt{instruccion}.\\
$A$: \nt{asignacion}; $E$: \nt{asignacion\_expresion}.\\
$X$: \nt{expresion}; $O$: \nt{operando}; $R$: \nt{operador}.

\begin{longtable}{L{1.3cm}L{\dimexpr\textwidth-1.3cm-4\tabcolsep\relax}}
\toprule
\textbf{Regla} & \textbf{Forma sentencial}\\\midrule
\endhead
Inicio & $I$\\
P06 & \nt{mientras\_instr}\\
P22 & \texttt{mientras} \nt{condicion} $L$ \texttt{fin\_mientras}\\
P23 & \texttt{mientras} $O$ \nt{operador\_rel} $O\ L$ \texttt{fin\_mientras}\\
P14 & \texttt{mientras x} \nt{operador\_rel} $O\ L$ \texttt{fin\_mientras}\\
P24 & \texttt{mientras x <} $O\ L$ \texttt{fin\_mientras}\\
P15 & \texttt{mientras x < 10} $L$ \texttt{fin\_mientras}\\
P02 & \texttt{mientras x < 10} $I$ \texttt{fin\_mientras}\\
P04 & \texttt{mientras x < 10} $A$ \texttt{fin\_mientras}\\
P08 & \texttt{mientras x < 10} $E$ \texttt{fin\_mientras}\\
P10 & \texttt{mientras x < 10 x =} $X$ \texttt{\# fin\_mientras}\\
P13 & \texttt{mientras x < 10 x =} $O\ R\ O$ \texttt{\# fin\_mientras}\\
P14 & \texttt{mientras x < 10 x = x} $R\ O$ \texttt{\# fin\_mientras}\\
P16 & \texttt{mientras x < 10 x = x +} $O$ \texttt{\# fin\_mientras}\\
P15 & \texttt{mientras x < 10 x = x + 1 \# fin\_mientras}\\
\bottomrule
\end{longtable}

\figuraamplia{figuras/arbol_extension.pdf}{e) Árbol sintáctico de la extensión creativa}{Árbol
de derivación de la instrucción de lazo \texttt{mientras x < 10 x = x + 1 \# fin\_mientras}.}{La
raíz \nt{instruccion} se expande con P06 hasta \nt{mientras\_instr}, que, a su vez, se expande con
P22. Se distinguen las cuatro partes de la construcción: la palabra reservada \texttt{mientras}, la
condición \nt{condicion}, el cuerpo \nt{lista\_instrucciones} y el cierre
\texttt{fin\_mientras}.}{fig:arbol-extension}

\section*{4. Conclusiones}

La gramática extendida describe el lenguaje del enunciado con 14 no terminales, 19 terminales y
27 producciones, con \nt{programa} como símbolo inicial. Tratar los identificadores y los números
como clases léxicas permite mantener un conjunto finito de reglas y separar el análisis léxico del
sintáctico. Las producciones distinguen la asignación simple, la asignación con expresión y la
salida. La validación acepta el programa de referencia completo y rechaza los cuatro casos
incorrectos del enunciado; cada cadena aceptada en las pruebas tiene un único árbol sintáctico.
La extensión \texttt{mientras ... fin\_mientras} incorpora repetición y anidamiento mediante nuevas
reglas que reutilizan los operandos y las listas de instrucciones de la gramática base.

JFLAP 7.1 trabaja con símbolos de un carácter, por lo que se definió una codificación isomorfa
para representar la misma gramática en el archivo entregado. Su biblioteca comprobó la lectura de
las 27 producciones y las capturas del anexo D documentan las tres derivaciones principales.
El analizador propio valida las cadenas por tokens y genera los árboles legibles. Estas pruebas
comprueban la estructura sintáctica de los casos ensayados; no comprueban declaración de variables,
tipos ni divisiones por cero, y tampoco constituyen una demostración general de ausencia de
ambigüedad.

\section*{5. Consideraciones de entrega}

La gramática para JFLAP se encuentra en GitHub mediante el
\href{https://github.com/FernandoValdes01/Lab03-TC/blob/main/gramatica/GLC_Lab03.cfg}{enlace directo a \texttt{GLC\_Lab03.cfg}}.
También se adjunta una
\href{https://github.com/FernandoValdes01/Lab03-TC/blob/main/gramatica/GLC_Lab03.jff}{copia idéntica con extensión \texttt{.jff}}.
El archivo usa la codificación de un carácter por símbolo descrita en el anexo D.

La carpeta de entrega contiene el informe y los siguientes archivos:

\begin{itemize}
  \item \texttt{Fernando\_Valdes\_Gomez.pdf}: informe del laboratorio.
  \item \texttt{gramatica/\allowbreak GLC\_Lab03.cfg}: gramática en el formato XML de JFLAP 7.1, con
  la codificación isomorfa de un carácter por símbolo (14 variables, 19 terminales, 27 producciones).
  \item \texttt{gramatica/\allowbreak GLC\_Lab03.jff}: copia idéntica del archivo anterior para abrir en JFLAP.
  \item \texttt{gramatica/\allowbreak GLC\_Lab03.txt}: definición formal legible de $G=(V,\Sigma,P,S)$
  y tabla de codificación para JFLAP.
  \item \texttt{herramientas/\allowbreak gramatica\_lab3.py}: especificación de la gramática (fuente única de
  verdad).
  \item \texttt{herramientas/\allowbreak validar\_gramatica.py}: analizador, generador de derivaciones y de
  árboles.
  \item \texttt{herramientas/\allowbreak generar\_cfg\_jflap.py}: generador del archivo de JFLAP y de
  la tabla de codificación.
  \item \texttt{herramientas/\allowbreak VerificarGramaticaJFLAP.\allowbreak java}: verificador del archivo de JFLAP con
  la biblioteca de JFLAP 7.1.
  \item \texttt{herramientas/\allowbreak CapturarArbolJFLAP.\allowbreak java}: harness que reproduce las capturas de los
  árboles con JFLAP 7.1.
  \item \texttt{herramientas/\allowbreak test\_validar\_gramatica.py}: pruebas de regresión del analizador y de la correspondencia con el XML.
  \item \texttt{herramientas/\allowbreak VerificarCasosJFLAP.java},
  \texttt{gramatica/\allowbreak casos\_jflap.tsv} y
  \texttt{evidencias/\allowbreak validacion\_casos\_jflap.txt}: verificación de 14 casos con la biblioteca de JFLAP.
  \item \texttt{gramatica/\allowbreak equivalencias\_jflap.txt}: correspondencia entre los símbolos y su codificación.
  \item \texttt{figuras/}: árboles sintácticos en PDF vectorial, con originales SVG y copias PNG, y
  capturas de JFLAP 7.1 en \texttt{jflap\_*.png}.
  \item \texttt{evidencias/\allowbreak evidencias\_lab03.txt} y \texttt{evidencias/\allowbreak verificacion\_jflap.txt}:
  registro de validación propia y de verificación del archivo de JFLAP, respectivamente.
  \item \texttt{README.md}: guía de la carpeta de entrega y de cómo reproducir las evidencias.
\end{itemize}

Las figuras~\ref{fig:arbol-zxy}, \ref{fig:arbol-salida} y \ref{fig:arbol-extension} son esquemas
legibles generados por el analizador propio; las figuras~\ref{fig:jflap-zxy},
\ref{fig:jflap-salida} y \ref{fig:jflap-extension} son capturas de JFLAP 7.1 de las mismas
cadenas. Las evidencias adicionales se reúnen en los anexos C, D y E.

\clearpage
\section*{Anexos}

\subsection*{Anexo A -- La gramática en su versión carácter a carácter}
La versión entregada trata los
identificadores y los números como clases léxicas. Si se quisiera llevar esa descripción al
interior de la gramática, bastaría reemplazar \texttt{IDENT} y \texttt{NUM} por reglas como
\nt{identificador} $\rightarrow$ \nt{letra} \nt{resto\_id},
\nt{resto\_id} $\rightarrow$ \nt{resto\_id} \nt{letra} \,|\, \nt{resto\_id} \nt{digito} \,|\,
\nt{resto\_id} \texttt{\_} \,|\, $\varepsilon$, y
\nt{numero} $\rightarrow$ \nt{entero} \,|\, \nt{entero} \texttt{.} \nt{entero},
\nt{entero} $\rightarrow$ \nt{digito} \nt{entero} \,|\, \nt{digito}. Esa versión exige enumerar las
52 letras y los 10 dígitos, lo que multiplica el tamaño de la gramática sin aportar capacidad
descriptiva: es la razón por la que se prefirió la versión con clases léxicas.

\subsection*{Anexo B -- Derivación completa del programa de referencia}
El archivo \texttt{evidencias/\allowbreak evidencias\_lab03.txt} contiene, en su sección 5,
las 108 aplicaciones de producciones.

\subsection*{Anexo C -- Registro de la validación automática}
El mismo archivo, en sus secciones 2, 2.b,
3, 6 y 7, registra el veredicto de los ocho casos, el número de árboles sintácticos de cada uno
(uno en todos los casos aceptados, lo que confirma análisis único en esos ejemplos), el análisis de las cadenas
adicionales y los resultados de la extensión creativa.

\apartado{Cadenas de control adicionales}

Para reforzar la validación se analizaron otras 24 cadenas, incluidos tres controles
Unicode; la tabla muestra las más representativas. Las rechazadas confirman, cada una, una
restricción del lenguaje.

\begin{longtable}{|L{6.6cm}|L{2.6cm}|L{\dimexpr\textwidth-6.6cm-2.6cm-6\tabcolsep-4\arrayrulewidth\relax}|}
\hline
\rowcolor{azulclaro}\textbf{Cadena} & \textbf{Resultado} & \textbf{Observación}\\\hline
\endfirsthead
\hline
\rowcolor{azulclaro}\textbf{Cadena} & \textbf{Resultado} & \textbf{Observación}\\\hline
\endhead

\texttt{x = 1.345 \#} & Aceptada & Número real con punto decimal.\\\hline
\texttt{m = k \% 10 \#} & Aceptada & Operador módulo (P20).\\\hline
\texttt{t = k/z\#} & Aceptada & División sin espacios (P19).\\\hline
\texttt{temporal = x \#} & Aceptada & Asignación del valor de otra variable (P12).\\\hline
\texttt{z = x + y * z \#} & Rechazada & Expresión con dos operadores: la gramática define la operación binaria del enunciado.\\\hline
\texttt{Begin Programa p End} & Rechazada & Cuerpo vacío: la lista de instrucciones exige al menos una (P02, P03).\\\hline
\texttt{x == 20 \#} & Rechazada & Doble signo igual: no existe la producción que lo permita.\\\hline
\texttt{9x = 20 \#} & Rechazada & Un identificador no puede empezar con un dígito.\\\hline
\texttt{salida 20 \#} & Rechazada & La salida exige un identificador, no un número.\\\hline
\texttt{SALIDA x \#} & Rechazada & Las palabras reservadas distinguen mayúsculas.\\\hline
\texttt{begin Programa p x = 1 \# End} & Rechazada & \texttt{Begin} es la única forma de abrir un programa.\\\hline
\end{longtable}

\subsection*{Anexo D -- Verificación y capturas de JFLAP 7.1}

JFLAP 7.1 no acepta terminales de varios caracteres: su analizador recorre el lado derecho de cada
producción carácter a carácter y considera variable a toda letra mayúscula. Para que el archivo
entregado represente la misma gramática sin cambiar el lenguaje, se usa una codificación isomorfa de
un carácter por símbolo. La tabla siguiente muestra el mapeo completo; la gramática formal del
informe sigue usando los nombres legibles.

\Needspace{10\baselineskip}
\begin{longtable}{|L{4.6cm}|L{1.4cm}|L{3.5cm}|L{1.4cm}|}
\hline
\rowcolor{azulclaro}\textbf{No terminal} & \textbf{Código} & \textbf{Terminal} & \textbf{Código}\\\hline
\endfirsthead
\hline
\rowcolor{azulclaro}\textbf{No terminal} & \textbf{Código} & \textbf{Terminal} & \textbf{Código}\\\hline
\endhead

\nt{programa} & \tm{S} & \texttt{Begin} & \tm{b}\\\hline
\nt{lista\_instrucciones} & \tm{L} & \texttt{End} & \tm{e}\\\hline
\nt{instruccion} & \tm{I} & \texttt{Programa} & \tm{p}\\\hline
\nt{asignacion} & \tm{A} & \texttt{salida} & \tm{s}\\\hline
\nt{asignacion\_simple} & \tm{B} & \texttt{mientras} & \tm{m}\\\hline
\nt{asignacion\_expresion} & \tm{E} & \texttt{fin\_mientras} & \tm{f}\\\hline
\nt{valor} & \tm{V} & \texttt{+} & \tm{+}\\\hline
\nt{expresion} & \tm{X} & \texttt{-} & \tm{-}\\\hline
\nt{operando} & \tm{O} & \texttt{*} & \tm{*}\\\hline
\nt{operador} & \tm{R} & \texttt{/} & \tm{/}\\\hline
\nt{salida\_instr} & \tm{Q} & \texttt{\%} & \tm{\%}\\\hline
\nt{mientras\_instr} & \tm{M} & \texttt{=} & \tm{=}\\\hline
\nt{condicion} & \tm{C} & \texttt{\#} & \tm{\#}\\\hline
\nt{operador\_rel} & \tm{T} & \texttt{<} & \tm{\textless}\\\hline
\multicolumn{4}{|l|}{\textbf{Terminales adicionales}}\\\hline
 & & \texttt{>} & \tm{\textgreater}\\\hline
 & & \texttt{<=} & \tm{l}\\\hline
 & & \texttt{>=} & \tm{g}\\\hline
 & & \texttt{IDENT} & \tm{i}\\\hline
 & & \texttt{NUM} & \tm{n}\\\hline
\end{longtable}

Con esa codificación, JFLAP lee 14 variables, 19 terminales y 27 producciones, y el símbolo inicial
es \texttt{S}, que corresponde a \nt{programa}. La salida del verificador está en
\texttt{evidencias/\allowbreak verificacion\_jflap.txt}. Las tres derivaciones se reproducen con las
cadenas codificadas de la tabla siguiente; en cada captura se indica el no terminal desde el que se
inicia la derivación, de modo que el árbol coincida con el esquema legible de las figuras~\ref{fig:arbol-zxy}, \ref{fig:arbol-salida} y \ref{fig:arbol-extension}.

\Needspace{8\baselineskip}
\begin{longtable}{|L{4.7cm}|L{3.1cm}|L{3.3cm}|}
\hline
\rowcolor{azulclaro}\textbf{Cadena legible} & \textbf{Cadena codificada} & \textbf{Inicio}\\\hline
\endfirsthead
\hline
\rowcolor{azulclaro}\textbf{Cadena legible} & \textbf{Cadena codificada} & \textbf{Inicio}\\\hline
\endhead

\texttt{z = x + y \#} & \tm{i=i+i\#} & \nt{instruccion}\\\hline
\texttt{salida z \#} & \tm{si\#} & \nt{instruccion}\\\hline
\texttt{mientras x < 10 x = x + 1 \# fin\_mientras} & \tm{mi\textless ni=i+n\#f} & \nt{mientras\_instr}\\\hline
\end{longtable}

\figuraamplia{figuras/jflap_arbol_z_xy.png}{a) Captura de JFLAP 7.1 para \texttt{z = x + y \#}}{Árbol
de derivación producido por JFLAP 7.1 al analizar la cadena codificada \texttt{i=i+i\#} desde
\nt{instruccion}.}{La raíz es \texttt{I}; las hojas \texttt{i}, \texttt{=}, \texttt{i}, \texttt{+},
\texttt{i} y \texttt{\#} corresponden a \texttt{z}, \texttt{=}, \texttt{x}, \texttt{+}, \texttt{y} y
\texttt{\#}. La barra inferior muestra el estado final de JFLAP: \texttt{Derivations complete}.}{fig:jflap-zxy}

\figuraamplia{figuras/jflap_arbol_salida_z.png}{b) Captura de JFLAP 7.1 para \texttt{salida z \#}}{Árbol
de derivación de la cadena codificada \texttt{si\#} desde \nt{instruccion}.}{La raíz es \texttt{I}, que
se expande con P05 hasta \texttt{Q}; las hojas \texttt{s}, \texttt{i} y \texttt{\#} son
\texttt{salida}, \texttt{z} y \texttt{\#}.}{fig:jflap-salida}

\figuraamplia{figuras/jflap_arbol_extension.png}{c) Captura de JFLAP 7.1 para la extensión creativa}{Árbol
de derivación de la cadena codificada \texttt{mi\textless ni=i+n\#f} desde \nt{mientras\_instr}.}{La
raíz \texttt{M} se expande con P22 y reproduce las cuatro partes del lazo: \texttt{mientras},
\nt{condicion}, \nt{lista\_instrucciones} y \texttt{fin\_mientras}. Esta captura ejercita P22,
P23 y P24; las producciones P25--P27 quedan definidas, pero el ejemplo no usa
\texttt{>}, \texttt{\textless=} ni \texttt{\textgreater=}.}{fig:jflap-extension}

\subsection*{Anexo E -- Árbol adicional de la instrucción de salida}

\figuraaccesible{figuras/arbol_salida_z.pdf}{Árbol sintáctico de \texttt{salida z \#}}{Árbol de
derivación de la cadena \texttt{salida z \#}.}{La raíz \nt{instruccion} se expande con P05 hasta
\nt{salida\_instr}, que a su vez se expande con P21. Las tres hojas son los terminales
\texttt{salida}, \texttt{z} y \texttt{\#}.}{fig:arbol-salida}

\end{document}
